\chapter{Análisis de las metodologías}

\begin{quotation}[Físico] {Isaac Newton (1643--1727)}
El método apropiado para obtener conocimiento sobre las propiedades de las cosas es deducirlas de experimentos.
\end{quotation}

\begin{abstract}
Antes de comenzar incluso con la planificación del proyecto se ha de definir cómo se hará, por ello analicemos las metodologías que se han tenido en cuento para el desarrollo del proyecto.
\end{abstract}

\section{Metodologías tradicionales frente a metodologías ágiles}

Para llevar a cabo el desarrollo de software es necesario que existan unas guías que regulen la forma de organizar el desarrollo del proyecto, así como su estructura y los fases que atraviesa un proyecto desde su inicio hasta su cierre.
Para dar solución a esta problemática, se definieron las primeras metodologías, conocidas como metodologías tradicionales. Estas metodologías sirvieron para definir una serie de fases: especificación de los requisitos del producto a desarrollar, así como el diseño, implementación, pruebas, integración y el mantenimiento del mismo. Es lo que se conocen como ciclo de vida predictivo o ciclo de vida en cascada, en los cuales hay que definir el alcance, el tiempo y el costo del proyecto en las fases tempranas del ciclo de vida.
No obstante, en la actualidad los clientes pueden no tener claro los requisitos finales de su proyecto, debido a la falta de entendimiento entre el cliente y el equipo encargado de desarrollar el producto, así como el desconocimiento de la tecnología. Esto provoca que surjan una serie de cambios a lo largo del proyecto y, lo que puede conllevar a tener que regresar a una fase del proyecto.
Las circunstancias en las que puede ser de utilidad emplear una metodología tradicional pueden ser:

\begin{itemize}

\item En proyectos en que el cliente debe tener claro desde el principio los requisitos que serán acordados de una vez, y para todo el proyecto, no se esperaran cambios en ellos. 

\item Si el equipo que va a desarrollar el producto tiene bastante experiencia y sabe estimarlo con garantías de poder cumplir con las expectativas del cliente.

\item El proyecto tiene un alcance limitado, definido en las fases tempranas del ciclo de vida del mismo.

\end{itemize}

Las metodologías tradicionales, promueven imponer una disciplina al proceso de desarrollo software y de esta manera volverlo predecible y que sea eficiente. Si queremos obtener los beneficios anteriormente citados es imprescindible una estimación rigurosa al principio del proyecto.

Las metodologías ágiles surgen como una alternativa a las metodologías tradicionales puesto que, nos ayudan a reducir la probabilidad de fracaso debido a la mala estimación en costos, tiempos y funcionalidades en un entorno cambiante como es el desarrollo de software.

Las circunstancias en las que puede ser de utilidad emplear una metodología ágil pueden ser:

\begin{itemize}

\item Es fundamental que el cliente sea participativo durante todo el desarrollo del proyecto, ya que al tratarse de un desarrollo nuevo, lo habitual es que los requisitos y las funcionalidades se deban ajustar y cambiar según lo acordado en cada momento.

\item Uno de los pilares de esta metodología es la cohesión del equipo que va a desarrollar el producto, y la capacidad de cada miembro del equipo de poder complementarse con el resto de miembros del equipo.

\item Los requisitos se mantienen como un documento vigente que se van actualizando habitualmente, y las prioridades del desarrollo se pueden cambiar conforme avanza el proyecto. Aunque la documentación en la metodología ágil no es una prioridad, es necesario disponer de una documentación corta y concisa que sea perfectamente exhaustiva.

\item En proyectos ágiles, el trabajo se divide en partes a las cuales se le asigna un nivel de priorización y se desarrollan durante un periodo de tiempo que puede durar entre una y cuatro semanas, denominados sprints. 

\item El objetivo es obtener un producto funcional lo antes posible, de esta manera podamos tener una primera versión del producto que se irá mejorando en con cada sprint y una ventaja de esto es poder enseñar al cliente una versión primitiva de su producto, lo que fomentará su grado de implicación y participación en el desarrollo del mismo. 

\item Dado que este tipo de metodologías incurren en mayores incertidumbre y riesgos. Para hacer frente, a esto, en cada iteración los riesgos serán identificados, analizados y gestionados durante la duración de la misma, a través de revisiones de los productos, a fin de garantizar que el riesgo sea comprendido y controlado.  

\end{itemize}

Las metodologías ágiles son adecuadas en el caso de que tener un entorno tecnológico muy cambiante donde no se tiene definido desde el principio los requisitos del proyecto y que pueden cambiar de forma previsible.

Las metodologías ágiles más utilizadas en las empresas actuales son: programación extrema (XP), Scrum y Kanban, las cuales se guían a través del patrón establecido en el Manifiesto Ágil.

\section{Scrum}

En cuanto a la metodología que se han seguido para el desarrollo del proyecto, se trata de una metodología ágil, puesto que de manera general durante el desarrollo de una aplicación surgen cambios a lo largo de del proyecto, por tanto la ventaja de las metodologías ágiles en cuanto a la realización de cambios en el proyecto durante el desarrollo de este, la hace ideal para el desarrollo. La metodología ágil que vamos a seguir será Scrum, ya que disponemos de una experiencia y conocimiento en esta metodología debido a que la hemos aplicado de manera recurrente a lo largo de nuestra vida académica.

La metodología de Scrum está enfocada para grupos relativamente grandes de trabajo por lo que para el equipo de desarrollo con una separación tan marcada como es nuestro equipo de desarrollo compuesto por dos desarrolladores hemos decidido adaptarla para adecuarla a nuestras necesidades. En primer lugar no realizábamos Sprint Dailys de manera tan estricta como determina Scrum, en vez de hacer reuniones diarias las realizábamos cada dos días. Organizábamos sesiones de trabajo en grupo que fomentasen la coordinación y cooperación por parte del equipo. Cada domingo se realizaba una retrospectiva de la semana para analizar el avance a mitad de iteración por si necesitara una replanificación así como el chequeo del avance de las tareas por si se necesitase en el caso de las tareas del Back-End una reasignación. Cada dos semanas se realiza una Sprint Retrospective para canalizar todo lo que haya ido mal en la iteración y mejorar de cara a la siguiente.  